% Самостійний документ. Компіляція: xelatex (двічі) або lualatex.
% Усі числові приклади та набори даних навчальні.
\documentclass[aspectratio=169,11pt]{beamer}
\usepackage{fontspec}
\usepackage{polyglossia}
\setdefaultlanguage{ukrainian}
\setmainfont{DejaVu Serif}
\setsansfont{DejaVu Sans}
\setmonofont{DejaVu Sans Mono}
\newfontfamily\cyrillicfont{DejaVu Serif}
\newfontfamily\cyrillicfontsf{DejaVu Sans}
\newfontfamily\cyrillicfonttt{DejaVu Sans Mono}
\usepackage{amsmath,amssymb,booktabs,array}
\usefonttheme{professionalfonts}
\definecolor{ink}{HTML}{183341}
\definecolor{accent}{HTML}{087C83}
\setbeamercolor{normal text}{fg=ink,bg=white}
\setbeamercolor{structure}{fg=accent}
\setbeamercolor{frametitle}{fg=ink}
\setbeamertemplate{navigation symbols}{}
\setbeamertemplate{footline}{\hfill\insertframenumber\hspace{5mm}\vspace{3mm}}
\setbeamersize{text margin left=8mm,text margin right=8mm}
\setbeamerfont{frametitle}{size=\large,series=\bfseries}
\hypersetup{colorlinks=true,urlcolor=accent,linkcolor=ink}
\newcommand{\unit}[1]{\,\text{#1}}
\newcommand{\src}[2]{\par\vfill{\tiny\href{#1}{#2}}}
\title{Засоби вимірювальної техніки\newline і вимірювальні канали}
\subtitle{Лекція 3. Математичний конспект і практичні завдання}
\author{Метрологія, технологічні вимірювання та прилади}
\date{}
\begin{document}
\begin{frame}
\titlepage
\end{frame}

\begin{frame}{Позначення та межі моделей}
\begin{itemize}
\item $x$ і $y$ позначають фізичні величини, $f$ їхній номінальний зв'язок.
\item $D$ є допустимою областю входу, $f(D)$ є множиною можливих виходів.
\item $p_{\mathrm{kPa}}=p/(1\unit{кПа})$, $I_{\mathrm{mA}}=I/(1\unit{мА})$ та $U_{\mathrm V}=U/(1\unit{В})$ є \textbf{числовими значеннями}.
\item Числовий коефіцієнт залежить від обраних одиниць. Фізична залежність від цього не змінюється.
\item Якщо не зазначено інше, режими статичні, характеристики лінійні, межі включені. Навчальні дані не є специфікацією реального виробу.
\end{itemize}
Невизначеності та похибки не підсумовуємо. Номінальна модель не доводить точності чи електричної безпеки каналу.
\end{frame}

\section{Класи ЗВТ}
\begin{frame}{Визначення: модель вимірювальної ланки}
Для аналізу інтерфейсів подамо ланку як запис
\[
L_i=(T_i^{\rm in},u_i^{\rm in},D_i,\;T_i^{\rm out},u_i^{\rm out},Y_i,\;f_i),
\qquad y_i=f_i(x_i),\quad x_i\in D_i.
\]
$T$ задає тип сигналу, $u$ одиницю, $D$ допустимі входи, $Y=f(D)$ можливі виходи. Для невідомої характеристики можна задати паспортний вихідний інтервал, який охоплює $f(D)$.
\medskip

Міра відтворює приписане значення. Перетворювач має визначений зв'язок входу й виходу. Прилад, установка та система відрізняються функціями й організацією вимірювання.
\medskip

Формальний запис ланки є навчальною моделлю інтерфейсу, а не універсальним означенням усіх ЗВТ.
\src{https://jcgm.bipm.org/vim/en/3.7.html}{VIM 3.7: measuring transducer}
\end{frame}

\begin{frame}{Визначення: сумісність сусідніх ланок}
За попередньо узгоджених одиниць необхідно
\[
T_i^{\rm out}=T_{i+1}^{\rm in},\qquad
u_i^{\rm out}=u_{i+1}^{\rm in},\qquad Y_i\subseteq D_{i+1}.
\]
Для замкнених інтервалів $Y_i=[a,b]$, $D_{i+1}=[c,d]$:
\[
[a,b]\subseteq[c,d]\quad\Longleftrightarrow\quad c\le a\le b\le d.
\]
\textbf{Приклад:} вихід $[1;5]\unit{В}$ сумісний за діапазоном із входом $[0;5]\unit{В}$, але не з $[0;3{,}3]\unit{В}$.
\medskip

Збіг одиниць не гарантує збігу фізичного змісту: абсолютний і надлишковий тиск потребують різних опор. Повна перевірка також враховує навантаження, живлення, підключення та швидкодію.
\end{frame}

\section{Функціональне і метрологічне призначення}
\begin{frame}{Визначення: лінійне масштабування}
Для відповідності $x_L\mapsto y_L$, $x_H\mapsto y_H$, де $x_H>x_L$:
\[
y=y_L+\frac{x-x_L}{x_H-x_L}(y_H-y_L)=ax+b,
\]
\[
a=\frac{y_H-y_L}{x_H-x_L},\qquad b=y_L-ax_L.
\]
$a$ має одиницю $[y]/[x]$, $b$ має одиницю виходу.
\medskip

Якщо $a\ne0$, обернене перетворення $x=(y-b)/a$ дозволене лише в області застосування моделі. Якщо $a=0$, різні входи дають однаковий вихід і таке відновлення неможливе.
\medskip

Ця залежність описує \textbf{функцію} перетворювача. Вона не встановлює його еталонного статусу або придатності для калібрування.
\end{frame}

\begin{frame}{Навчальний приклад: струм у напругу}
Ідеалізоване приймання струму на резисторі:
\[
U=RI,\qquad R=250\unit{Ом},\qquad I\in[0{,}004;0{,}020]\unit{А}.
\]
Тоді $U\in[1;5]\unit{В}$. У числових значеннях:
\[
U_{\mathrm V}=0{,}25I_{\mathrm{mA}}.
\]
\begin{center}
\begin{tabular}{c|ccc}
$I$, мА & 4 & 12 & 20\\\hline
$U$, В & 1 & 3 & 5
\end{tabular}
\end{center}
Для входу $U\le3{,}3\unit{В}$ допустимо лише $I\le13{,}2\unit{мА}$.
\medskip

Припущення: струмова петля підтримує струм, вхід майже не навантажує резистор. Допустима потужність і запас живлення потребують окремої перевірки.
\end{frame}

\begin{frame}{Визначення: область придатного перетворення}
Для вхідного діапазону $D$ і допустимого виходу приймача $E$:
\[
D_{\rm valid}=\{x\in D:f(x)\in E\}=D\cap f^{-1}(E).
\]
Тут $f^{-1}(E)$ є \textbf{прообразом множини}, а не обов'язково оберненою функцією.
\medskip

Якщо $f(x)=ax+b$, $a>0$ та $E=[y_L,y_H]$, то
\[
D_{\rm valid}=D\cap\left[\frac{y_L-b}{a},\frac{y_H-b}{a}\right].
\]
Для $a<0$ межі міняються місцями. За $a=0$ допустимий весь $D$, якщо $b\in E$, і жодного входу інакше.
\medskip

Нормовану характеристику застосовують лише за встановлених умов. Межа вимірювання та межа допустимого перевантаження не тотожні.
\end{frame}

\section{Еталони і передавання розміру одиниці}
\begin{frame}{Визначення: опорне значення та еталон}
Еталон забезпечує опорне значення величини з пов'язаною невизначеністю. Скорочений запис:
\[
\text{опора}=(q_{\rm ref},\;u(q_{\rm ref}),\;\text{умови та документований зв'язок}).
\]
Це схема зберігання відомостей, а не формула повного результату.
\begin{itemize}
\item Національний еталон має офіційне визнання як основа для приписування значень іншим еталонам відповідної величини.
\item Вторинний еталон установлюють калібруванням відносно первинного.
\item Робочий еталон регулярно використовують для інших приладів або систем.
\end{itemize}
Класифікаційні ознаки можуть поєднуватися. Робочий еталон і робочий датчик мають різні ролі.
\src{https://jcgm.bipm.org/vim/en/5.1.html}{VIM 5.1; також 5.3, 5.5, 5.7}
\end{frame}

\begin{frame}{Навчальна модель опорного тиску}
Для поршня в механічній рівновазі спрощено:
\[
p\approx\frac{F}{A_{\rm eff}}\approx\frac{mg}{A_{\rm eff}}.
\]
$m$ позначає масу навантаження, $g$ місцеве прискорення вільного падіння, $A_{\rm eff}$ ефективну площу поршневої пари.
\[
\frac{\unit{кг}\cdot\unit{м}/\unit{с}^{2}}{\unit{м}^{2}}
=\unit{Н}/\unit{м}^{2}=\unit{Па}.
\]
Реальна модель враховує поправки та умови. Ефективну площу не замінюють грубою геометричною оцінкою.
\medskip

NIST описує передавання характеристики через зрівноважування тисків двох поршневих манометрів. Наведене рівняння пояснює фізичну основу, але не відтворює повної процедури NIST.
\src{https://www.nist.gov/news-events/news/2017/05/faster-way-calibrate-piston-gauges}{NIST: A Faster Way to Calibrate Piston Gauges, 2017}
\end{frame}

\begin{frame}{Визначення: поправка в одній робочій точці}
Опорний засіб має показ $r_i$ і відому поправку $c_i$. Наступний засіб при тому самому вході має показ $d_i$:
\[
q_i=r_i+c_i,\qquad c_{i+1}=q_i-d_i.
\]
Отже, $d_i+c_{i+1}=q_i$ в номінальній моделі цієї точки.
\medskip

\textbf{Окремий приклад:} $r_i=100{,}00\unit{кПа}$, $c_i=+0{,}02\unit{кПа}$, $d_i=100{,}08\unit{кПа}$.
\[
q_i=100{,}02\unit{кПа},\qquad c_{i+1}=-0{,}06\unit{кПа}.
\]
Поправку додають до показу. Її не переносять на інші точки без моделі залежності. Рівність після поправки не означає нульової невизначеності.
\src{https://jcgm.bipm.org/vim/en/2.53.html}{VIM 2.53: correction}
\end{frame}

\begin{frame}{Послідовність порівнянь і простежуваність}
У спрощеній моделі однієї точки:
\[
c_{i+1}=r_i+c_i-d_i,\qquad i=0,\ldots,n-1.
\]
Кожний крок має власні покази, умови та запис порівняння. Наступний опорний засіб є попереднім порівнюваним засобом.
\medskip

Числова рекурентність сама не доводить метрологічної простежуваності. Для неї потрібен неперервний документований калібрувальний зв'язок результату з опорою, з урахуванням внесків у невизначеність.
\medskip

Невизначеності зберігаємо як дані кроків. Не додаємо їх арифметично та не застосовуємо корінь із суми квадратів без моделі й припущень.
\src{https://jcgm.bipm.org/vim/en/2.41.html}{VIM 2.41: metrological traceability}
\end{frame}

\section{Структура каналу і базові метрологічні характеристики}
\begin{frame}{Визначення: композиція каналу}
Нехай $x_0=x$, а $x_i=f_i(x_{i-1})$. Тоді
\[
y=F(x)=(f_n\circ\cdots\circ f_1)(x).
\]
Композицію читаємо справа наліво: спочатку $f_1$, потім $f_2$ і так далі.
\[
D_{\rm channel}=\{x\in D_1: x_{i-1}\in D_i\ \text{для всіх }i=2,\ldots,n\}.
\]
За потреби додатково вимагаємо $y\in E$ для приймача результату.
\medskip

Для $f_i(z)=a_i z+b_i$ та двох ланок:
\[
F(x)=a_2a_1x+(a_2b_1+b_2).
\]
Порядок важливий: зазвичай $f_2\circ f_1\ne f_1\circ f_2$, а фізично переставлена послідовність може бути взагалі несумісною.
\end{frame}

\begin{frame}{Визначення: чутливість}
Для лінійної характеристики $y=ax+b$:
\[
S=\frac{\Delta y}{\Delta x}=a\qquad(\Delta x\ne0).
\]
Для диференційовної нелінійної характеристики в точці $x$:
\[
S(x)=\frac{df}{dx}(x).
\]
Чутливість є нахилом, її одиниця $[y]/[x]$. Для послідовності диференційовних ланок правило ланцюга дає
\[
\frac{dF}{dx}=\prod_{i=1}^{n} f_i'(x_{i-1}).
\]
Для лінійних ланок це $S_{\rm total}=\prod_i a_i$. Формула застосовна всередині придатної області, без насичення та квантування.
\src{https://jcgm.bipm.org/vim/en/4.12.html}{VIM 4.12: sensitivity of a measuring system}
\end{frame}

\begin{frame}{Навчальний канал тиску}
У числових значеннях для $0\le p_{\mathrm{kPa}}\le1000$:
\[
I_{\mathrm{mA}}=4+0{,}016p_{\mathrm{kPa}},\qquad
U_{\mathrm V}=0{,}25I_{\mathrm{mA}}=1+0{,}004p_{\mathrm{kPa}}.
\]
Чутливості в одиницях фізичних величин:
\[
S_1=0{,}016\frac{\unit{мА}}{\unit{кПа}},\quad
S_2=0{,}25\frac{\unit{В}}{\unit{мА}},\quad
S_{\rm total}=0{,}004\frac{\unit{В}}{\unit{кПа}}.
\]
\textbf{Окрема ілюстративна точка:} $p=400\unit{кПа}$ дає $I=10{,}4\unit{мА}$ і $U=2{,}6\unit{В}$.
\medskip

Для $U\le3{,}3\unit{В}$ маємо $p\le575\unit{кПа}$. Весь придатний інтервал цього каналу: $[0;575]\unit{кПа}$.
\end{frame}

\begin{frame}{Визначення: поріг і варіація}
Поріг розрізнення описує найбільшу зміну входу, яка ще не викликає помітної зміни показу за визначених умов. Критерій помітності потрібно встановити окремо.
\medskip

Варіація при однаковому вході $x$ у прийнятій моделі:
\[
h(x)=\left|y_{\uparrow}(x)-y_{\downarrow}(x)\right|.
\]
Індекси означають підхід до точки при збільшенні та зменшенні входу. Інші умови однакові, покази усталені.
\medskip

При $p=500\unit{кПа}$ навчальні покази $12{,}1$ і $11{,}9\unit{мА}$ дають $h=0{,}2\unit{мА}$. Це різниця гілок, не повна похибка та не стандартне відхилення.
\src{https://jcgm.bipm.org/vim/en/4.16.html}{VIM 4.16: discrimination threshold; також VIM 4.14: resolution}
\end{frame}

\begin{frame}{Визначення: час установлення}
За стрибка входу в момент $t=0$ і кінцевого виходу $y_\infty$:
\[
t_s=\inf\{t\ge0:\ |y(\tau)-y_\infty|\le\varepsilon\ \text{для всіх }\tau\ge t\}.
\]
$\varepsilon>0$ є заданою півшириною смуги. Якщо умова ніколи не виконується, скінченного $t_s$ немає.
\medskip

Навчальний запис: $t=[0;0{,}1;0{,}2;0{,}3;0{,}4;0{,}5;0{,}6]\unit{с}$,
\[
y=[0;0{,}70;0{,}98;1{,}10;1{,}03;1{,}01;1{,}00],\quad
y_\infty=1,\quad\varepsilon=0{,}05.
\]
У записі точки після $0{,}4\unit{с}$ залишаються в смузі. Це лише оцінка за наявними відліками, не доказ поведінки між ними або після завершення запису.
\src{https://jcgm.bipm.org/vim/en/4.23.html}{VIM 4.23: step response time}
\end{frame}

\begin{frame}{Джерела та відповідність курсу}
\small
Тема і чотири розділи: \textbf{Master\_Course\_Plan\_Metrology\_18\_lectures.pdf}, с. 12--13. Екзаменаційні питання 14, 15, 16 відтворені в розділі 8 плану, с. 65. Окремий DOCX недоступний.
\medskip

Термінологія: \href{https://jcgm.bipm.org/vim/en/}{JCGM, International Vocabulary of Metrology (VIM3)}, пункти 2.41, 2.53, 3.1--3.10, 4.7, 4.12, 4.14, 4.16, 4.23, 5.1--5.7.
\medskip

Реальні контексти:
\begin{itemize}
\item \href{https://www.ni.com/en/shop/data-acquisition/sensor-fundamentals/measuring-pressure-with-bridge-based-and-other-pressure-sensors.html}{NI: Measuring Pressure with Bridge-Based and Other Pressure Sensors}.
\item \href{https://www.ni.com/docs/en-US/bundle/ni-daqmx/page/sigcon.html}{NI: Signal Conditioning}.
\item \href{https://www.fluke.com/en-gb/product/calibration-tools/multifunction-calibrators/fluke-754}{Fluke: 754 Documenting Process Calibrator-HART}.
\item \href{https://www.nist.gov/news-events/news/2017/05/faster-way-calibrate-piston-gauges}{NIST: A Faster Way to Calibrate Piston Gauges}, 2017.
\end{itemize}
Числові моделі й завдання авторські навчальні, без готових рішень Python.
\end{frame}

\appendix
\begin{frame}[t]{Практичне завдання 1. Перевірка стиків}
\small
\textbf{Мета:} перевірити сумісність опису сусідніх ланок. Подайте кожну словником із назвою, типом, одиницею та межами входу й виходу.
\begin{center}\begin{tabular}{lll}\toprule
Ланка & Вхід & Вихід\\\midrule
Передавач & тиск, кПа, $[0;1000]$ & струм, мА, $[4;20]$\\
Перетворювач & струм, мА, $[4;20]$ & напруга, В, $[1;5]$\\
Контролер & напруга, В, $[0;5]$ & код, 1, $[0;4095]$\\\bottomrule
\end{tabular}\end{center}
\textbf{Завдання:} перевірте наявність полів, правильний порядок меж, типи й одиниці сусідів, повне включення вихідного інтервалу у вхідний. Автоматично одиниці не перераховуйте. Межі включені.
\par\medskip
\textbf{Тести:} базовий список; заміна входу контролера на струм у мА; звуження його входу до $[0;3{,}3]\unit{В}$; відсутня верхня межа. Зміни робіть у незалежних копіях.
\par\medskip
\textbf{Результат:} словник \texttt{compatible} і \texttt{reasons} із назвами проблемних ланок. Додайте коротке пояснення, чому позитивний висновок не перевіряє живлення й навантаження. Лише списки, словники, цикл та умови.
\end{frame}

\begin{frame}[t]{Практичне завдання 2. Масштаб і статус}
\small
\textbf{Мета:} описати приймання струму через $R=250\unit{Ом}$ функцією \texttt{convert(current\_mA, low\_V, high\_V)}.
\par\medskip
\textbf{Вхідні дані:} струми $[0;4;8;12;13{,}2;16;20;21]\unit{мА}$, два входи контролера $[0;5]\unit{В}$ та $[0;3{,}3]\unit{В}$.
\par\medskip
\textbf{Правила:} допустимий струм $[4;20]\unit{мА}$; усі межі включені. Використайте фізичну модель $U=RI$ з правильними одиницями. Некоректні межі або відсутній аргумент дають \texttt{invalid\_config}. Струм поза діапазоном дає \texttt{invalid\_current} і \texttt{voltage\_V=None}.
\par\medskip
Для дозволеного струму збережіть обчислену напругу. Якщо вона поза входом, статус \texttt{input\_overrange}, інакше \texttt{ok}. Окремо перевірте обидві точні границі струму та напруги.
\par\medskip
\textbf{Результат:} список словників \texttt{current\_mA}, \texttt{voltage\_V}, \texttt{status}; пояснення відмінностей двох конфігурацій. Не моделюйте діагностичні режими конкретних виробів і не обрізайте напругу до межі. Готову реалізацію не наведено.
\end{frame}

\begin{frame}[t]{Практичне завдання 3. Журнал передавання}
\small
\textbf{Мета:} передати поправки в одній точці. Початково $c_0=+0{,}03\unit{кПа}$. Усі покази в кПа. Умови незмінні, поправки сталі в цій точці.
\begin{center}\begin{tabular}{ccccc}\toprule
$i$ & $r_i$ & $d_i$ & документ & задане $u_i$, кПа\\\midrule
0 & 100,00 & 100,11 & C-01 & 0,01\\
1 & 100,09 & 99,96 & C-02 & 0,02\\
2 & 99,98 & 100,05 & C-03 & 0,03\\\bottomrule
\end{tabular}\end{center}
\textbf{Завдання:} використайте означення виправленої опори й поправки з розділу 3. На кожному кроці збережіть $q_i$, нову $c_{i+1}$, одиницю, документ та задане $u_i$. Невизначеності не підсумовуйте.
\par\medskip
\textbf{Тести:} повний список; копія з порожнім C-02; копія з відсутнім $d_1$. Числову обробку зупиніть там, де для неї бракує даних. Відсутній документ позначте окремо.
\par\medskip
\textbf{Результат:} журнал кроків, статус повноти й список проблем. Поясніть знак поправки та чому повний список не доводить простежуваності.
\end{frame}

\begin{frame}[t]{Практичне завдання 4. Модель каналу}
\small
\textbf{Мета:} побудувати окремі функції $p\mapsto I$ та $I\mapsto U$ і функцію проходження всього каналу. Параметри: $p\in[0;1000]\unit{кПа}$, $I\in[4;20]\unit{мА}$, $R=250\unit{Ом}$.
\par\medskip
\textbf{Тести:} $p=[-1;0;250;575;600;1000;1001]\unit{кПа}$ для входів контролера $[0;5]\unit{В}$ та $[0;3{,}3]\unit{В}$. Межі включені. Статуси: \texttt{ok}, \texttt{sensor\_overrange}, \texttt{input\_overrange}.
\par\medskip
\textbf{Завдання:} повертайте журнал із $p$, $I$, $U$ та причиною непридатності. За недопустимого $p$ не формуйте нормальний результат наступних ланок. За перевищення входу збережіть $U$ як діагностичне значення. Не обрізайте сигнал.
\par\medskip
Обчисліть добуток чутливостей із правильними одиницями. Перевірте його відношенням приростів для двох близьких допустимих тисків, обраних самостійно. АЦП та його квантування не моделюйте.
\par\medskip
\textbf{Результат:} таблиця проходження для обох конфігурацій, чутливість і короткий висновок про ланку-обмежувач. Не обчислюйте сумарної похибки. Потрібні лише функції, цикл, словники й умови.
\end{frame}
\end{document}
